% !TEX root = ../main.tex
% Responsible: C -- Reviewer: B   (see homework2/README.md, Division of work)
\subsection{Factor models (16--19, 21)}
Each stock is mapped on the five factors by OLS on the window, $r_{i,t} = a_i + \bm b_i'\bm f_t + e_{i,t}$; the EWP
has $\bm b_{EWP} = \frac1{16}\sum_i \bm b_i$. Each factor follows the AR(1)-GJR-GARCH of~\eqref{eq:garch} and the
correlation of the standardised factor residuals is either constant, CCC \citep{bollerslev1990}, or DCC(1,1)
\citep{engle2002}. With $\bm\Sigma_{f,t}$ the predicted factor covariance:
\begin{itemize}
  \item \textbf{Gaussian (16--17)}: $\VaR_{i,t}^\alpha = a_i + \bm b_i'\bm\mu_{f,t} + z_\alpha\sqrt{\bm b_i'\bm\Sigma_{f,t}\bm b_i + s^2_{e,i}}$,
        including the idiosyncratic variance $s^2_{e,i}$ (for the EWP, $\bm w'\widehat{\mathrm{Cov}}(\bm e)\bm w$).
  \item \textbf{FHS (18--19)}: the factor innovations are devolatilised and decorrelated with the Cholesky factor of
        $\bm R_s$, resampled together with the mapping residuals of the same day, and recombined with the
        volatilities and correlation predicted for day $t$; the VaR is the empirical quantile of 10{,}000 simulated returns.
  \item \textbf{Model 21}: $X_{i,t} = \bm b_i'\bm f_t$ is modelled with an AR(1)-GJR-GARCH(1,1) with skewed-$t$
        innovations \citep{hansen1994}, and the VaR of $X_{i,t}$ is used as VaR of the stock. Since $X$ contains no
        idiosyncratic risk, we expect too many violations for single stocks.
\end{itemize}
\todo{choice to discuss: factors are in USD and returns in CHF.}

\paragraph{Bonus models (22+, optional).} \todo{member C: additional models, if any (code in var\_factors.py).}
